\documentclass[12pt,a4paper]{article}
\usepackage[utf8]{inputenc}
\usepackage[T1]{fontenc}
\usepackage[french]{babel}
\usepackage{geometry}
\usepackage{graphicx}
\usepackage{listings}
\usepackage{xcolor}
\usepackage{booktabs}
\usepackage{array}
\usepackage{longtable}
\usepackage{hyperref}
\usepackage{fancyhdr}
\usepackage{titlesec}

% Configuration de la page
\geometry{margin=2.5cm}

% Configuration des en-têtes et pieds de page
\pagestyle{fancy}
\fancyhf{}
\fancyhead[L]{TP Slurm 1}
\fancyhead[R]{Gestion de cluster}
\fancyfoot[C]{\thepage}

% Configuration des listings pour le code bash
\definecolor{codebg}{RGB}{245,245,245}
\definecolor{codeframe}{RGB}{200,200,200}
\definecolor{codecomment}{RGB}{0,128,0}
\definecolor{codekeyword}{RGB}{0,0,255}
\definecolor{codestring}{RGB}{163,21,21}

\lstdefinestyle{bash}{
    backgroundcolor=\color{codebg},
    frame=single,
    framerule=0.5pt,
    rulecolor=\color{codeframe},
    basicstyle=\ttfamily\small,
    breaklines=true,
    breakatwhitespace=true,
    columns=fullflexible,
    keepspaces=true,
    showstringspaces=false,
    commentstyle=\color{codecomment},
    keywordstyle=\color{codekeyword},
    stringstyle=\color{codestring},
    tabsize=4,
    xleftmargin=0.5cm,
    xrightmargin=0.5cm,
}

\lstset{style=bash}

% Configuration des titres
\titleformat{\section}{\Large\bfseries}{\thesection}{1em}{}
\titleformat{\subsection}{\large\bfseries}{\thesubsection}{1em}{}
\titleformat{\subsubsection}{\normalsize\bfseries}{\thesubsubsection}{1em}{}

% Commandes personnalisées
\newcommand{\code}[1]{\texttt{#1}}
\newcommand{\command}[1]{\texttt{\$ #1}}

\begin{document}

% Page de titre
\begin{titlepage}
    \centering
    \vspace*{2cm}
    {\Huge\bfseries Rapport de TP Slurm\par}
    \vspace{1cm}
    {\Large Gestion de cluster et ordonnancement de jobs\par}
    \vspace{2cm}
    {\large Yahdhih Abdelwedoud\par}
    \vspace{1cm}
    {\large 5 février 2026\par}
    \vfill
    {\large ENSIIE - Semestre 4\par}
\end{titlepage}

\tableofcontents
\newpage

%%%%%%%%%%%%%%%%%%%%%%%%%%%%%%%%%%%%%%%%%%%%%%%%%%%%%%%%%%%%%%%%%%%%%%%%%%%%%%%
\section{Préparation et validation de l'environnement}
%%%%%%%%%%%%%%%%%%%%%%%%%%%%%%%%%%%%%%%%%%%%%%%%%%%%%%%%%%%%%%%%%%%%%%%%%%%%%%%

\subsection{Q1 : Démarrage des machines virtuelles}

\textbf{Objectif :} Démarrer 4 machines virtuelles avec PCOCC en utilisant le template du TP Slurm 1.

\textbf{Commande exécutée :}
\begin{lstlisting}
pcocc alloc -c2 mycentos74-tp-slurm1:4
\end{lstlisting}

\textbf{Résultat obtenu :}
\begin{lstlisting}
-bash-4.2$ pcocc alloc -c2 mycentos74-tp-slurm1:4
salloc: Granted job allocation 37993
Configuring hosts... (done)
\end{lstlisting}

\textbf{Analyse :} Le message confirme le succès du lancement des 4 machines virtuelles avec l'allocation du job 37993.

\subsection{Q2 : Suivi du démarrage de la première VM}

\textbf{Objectif :} Observer l'avancement du démarrage de la première machine virtuelle.

\textbf{Commande exécutée :}
\begin{lstlisting}
pcocc console
\end{lstlisting}

\textbf{Résultat obtenu :}
\begin{lstlisting}
(pcocc/37993) bash-4.2$ pcocc console
Warning: Permanently added 'hpc06,192.168.1.26' (ECDSA) to the list of known hosts.

Detaching ...
Killed by signal 15.
\end{lstlisting}

\textbf{Analyse :} La commande a été lancée tardivement, ce qui n'a pas permis de capturer l'état de démarrage des machines virtuelles en temps réel.

\subsection{Q3 : Sortie de la console avec Ctrl+C}

\textbf{Objectif :} Sortir de la commande \code{pcocc console} en utilisant Ctrl+C trois fois de suite.

\textbf{Résultat obtenu :}
\begin{lstlisting}
(pcocc/37993) bash-4.2$ pcocc console
Warning: Permanently added 'hpc06,192.168.1.26' (ECDSA) to the list of known hosts.

Detaching ...
Killed by signal 15.
(pcocc/37993) bash-4.2$ pcocc ssh vm0
[yahdhih.abdelwedoud@vm0 ~]$ sinfo -l
\end{lstlisting}

\textbf{Analyse :} Après la sortie de la console, la connexion SSH à vm0 a été établie avec succès.

\subsection{Q4 : Vérification de l'état du cluster}

\textbf{Objectif :} Se connecter à la machine virtuelle vm0 et vérifier l'état du cluster.

\textbf{Commandes exécutées :}
\begin{lstlisting}
pcocc ssh vm0
sinfo -l
\end{lstlisting}

\textbf{Résultat obtenu :}
\begin{lstlisting}
Mon Feb  2 15:46:05 2026
PARTITION AVAIL  TIMELIMIT   JOB_SIZE ROOT  SHARE  GROUPS  NODES  STATE NODELIST
vm*          up   infinite 1-infinite   no     NO     all      4   idle vm[0-3]
\end{lstlisting}

\textbf{Analyse :} Le cluster est opérationnel avec 4 nœuds disponibles.

\subsection{Q5 : Description du cluster}

D'après la sortie de la commande \code{sinfo -l}, voici les caractéristiques du cluster :

\begin{table}[h]
\centering
\begin{tabular}{|l|l|p{8cm}|}
\hline
\textbf{Paramètre} & \textbf{Valeur} & \textbf{Description} \\
\hline
PARTITION & vm* & Partition par défaut (marquée *) \\
AVAIL & up & Partition disponible \\
TIMELIMIT & infinite & Pas de limite de temps d'exécution \\
JOB\_SIZE & 1-infinite & Jobs de 1 à N tâches autorisés \\
ROOT & no & Exécution en root non autorisée \\
SHARE & NO & Allocation exclusive des ressources \\
GROUPS & all & Tous les groupes d'utilisateurs autorisés \\
NODES & 4 & Nombre de nœuds disponibles \\
STATE & idle & Tous les nœuds sont libres \\
NODELIST & vm[0-3] & Liste des nœuds : vm0, vm1, vm2, vm3 \\
\hline
\end{tabular}
\caption{Caractéristiques du cluster Slurm}
\end{table}

\subsection{Q6 : Première exécution interactive}

\textbf{Objectif :} Vérifier le bon fonctionnement de Slurm avec une exécution interactive.

\textbf{Première tentative :}
\begin{lstlisting}
srun -n 4 -N 4 hostname
\end{lstlisting}

\textbf{Erreur obtenue :}
\begin{lstlisting}
srun: error: Unable to allocate resources: Invalid account or account/partition combination specified
\end{lstlisting}

\textbf{Résolution :} Ajout de l'utilisateur dans la base de données de comptabilité Slurm :
\begin{lstlisting}
sudo sacctmgr -i add user name=yahdhih.abdelwedoud cluster=ensiie account=guests fairshare=10 qos=normal,debug
\end{lstlisting}

\textbf{Résultat :}
\begin{lstlisting}
 Adding User(s)
  yahdhih.abdelwedoud
 Settings =
  Default Account = (null)
 Associations =
  U = yahdhih.a A = guests     C = ensiie
 Non Default Settings
  Fairshare     = 10
  QOS           = debug,normal
\end{lstlisting}

\textbf{Analyse :} L'utilisateur a été ajouté avec succès au compte \code{guests} avec les QoS \code{normal} et \code{debug}.

%%%%%%%%%%%%%%%%%%%%%%%%%%%%%%%%%%%%%%%%%%%%%%%%%%%%%%%%%%%%%%%%%%%%%%%%%%%%%%%
\section{Premier script batch}
%%%%%%%%%%%%%%%%%%%%%%%%%%%%%%%%%%%%%%%%%%%%%%%%%%%%%%%%%%%%%%%%%%%%%%%%%%%%%%%

\subsection{Q1 : Créer un premier script affichant le nom du nœud}

Le script \code{script-1.sh} contient la commande \code{hostname} et retourne un code 0.

\begin{lstlisting}
#!/bin/bash
hostname
exit 0
\end{lstlisting}

Exécution locale (sans Slurm) :
\begin{lstlisting}
[yahdhih.abdelwedoud@vm0 ~]$ chmod +x script-1.sh
[yahdhih.abdelwedoud@vm0 ~]$ ./script-1.sh
vm0
\end{lstlisting}

\subsection{Q2 : Exécuter avec sbatch et analyser avec sacct}

\begin{lstlisting}
[yahdhih.abdelwedoud@vm0 ~]$ sbatch ./script-1.sh
Submitted batch job 57
[yahdhih.abdelwedoud@vm0 ~]$ sacct -j 57
    JobID    JobName  Partition    Account  AllocCPUS      State ExitCode
------------ ---------- ---------- ---------- ---------- ---------- --------
57           script-1.+         vm     guests          1  COMPLETED      0:0
57.batch          batch                guests          1  COMPLETED      0:0
\end{lstlisting}

Le job 57 s'est exécuté avec succès (State = COMPLETED, ExitCode = 0:0).

\subsection{Q3 : Modifier le script pour obtenir un état FAILED}

Pour obtenir un état FAILED, il suffit de retourner un code différent de 0 (par exemple \code{exit 1}).

\begin{lstlisting}
[yahdhih.abdelwedoud@vm0 ~]$ sbatch ./script-1.sh
Submitted batch job 60
[yahdhih.abdelwedoud@vm0 ~]$ sacct -j 60
    JobID    JobName  Partition    Account  AllocCPUS      State ExitCode
------------ ---------- ---------- ---------- ---------- ---------- --------
60           script-1.+         vm     guests          1     FAILED      1:0
60.batch          batch                guests          1     FAILED      1:0
\end{lstlisting}

Le job 60 a échoué avec ExitCode = 1:0.

\subsection{Q4 : Demander l'allocation de 4 tâches}

Après avoir remis le script avec \code{exit 0}, on demande 4 tâches avec \code{-n 4} :

\begin{lstlisting}
[yahdhih.abdelwedoud@vm0 ~]$ sbatch -n 4 script-1.sh
Submitted batch job 61
[yahdhih.abdelwedoud@vm0 ~]$ sacct -j 61
    JobID    JobName  Partition    Account  AllocCPUS      State ExitCode
------------ ---------- ---------- ---------- ---------- ---------- --------
61           script-1.+         vm     guests          4  COMPLETED      0:0
61.batch          batch                guests          2  COMPLETED      0:0
\end{lstlisting}

\textbf{Différence au niveau de l'accounting :}

\begin{table}[h]
\centering
\begin{tabular}{|l|c|c|}
\hline
\textbf{Job} & \textbf{AllocCPUS (job)} & \textbf{AllocCPUS (batch)} \\
\hline
57 (1 tâche) & 1 & 1 \\
61 (4 tâches) & 4 & 2 \\
\hline
\end{tabular}
\caption{Comparaison de l'allocation CPU}
\end{table}

Avec \code{-n 4}, Slurm alloue \textbf{4 CPUs} au job global, mais le script batch lui-même n'utilise que \textbf{2 CPUs} (sur un seul nœud). Les 4 CPUs sont réservés pour d'éventuelles sous-tâches (srun).

%%%%%%%%%%%%%%%%%%%%%%%%%%%%%%%%%%%%%%%%%%%%%%%%%%%%%%%%%%%%%%%%%%%%%%%%%%%%%%%
\section{Intégration de « step » parallèles}
%%%%%%%%%%%%%%%%%%%%%%%%%%%%%%%%%%%%%%%%%%%%%%%%%%%%%%%%%%%%%%%%%%%%%%%%%%%%%%%

L'exécution d'un script batch demandant l'allocation de plusieurs cœurs nécessite le lancement d'un « jobstep » afin de pouvoir les utiliser en parallèle. Le lancement d'un step se fait par l'intermédiaire d'une commande \code{srun} intégrée dans le script batch.

\subsection{Q1 : Intégrer l'exécution d'un step affichant le nom des machines}

\textbf{Script \code{sript-1-2.sh} :}
\begin{lstlisting}
#!/bin/bash
srun hostname
exit 0
\end{lstlisting}

\textbf{Exécution avec sbatch :}
\begin{lstlisting}
[yahdhih.abdelwedoud@vm0 ~]$ sbatch -n 4 sript-1-2.sh
Submitted batch job 58
\end{lstlisting}

\textbf{Contenu du fichier de sortie \code{slurm-58.out} :}
\begin{lstlisting}
vm0
vm0
vm1
vm1
\end{lstlisting}

\textbf{Données d'accounting :}
\begin{lstlisting}
[yahdhih.abdelwedoud@vm0 ~]$ sacct -j 58 -o jobid,jobname,nodelist,alloccpus,state
       JobID    JobName        NodeList  AllocCPUS      State 
------------ ---------- --------------- ---------- ---------- 
58           sript-1-2+         vm[0-1]          4  COMPLETED 
58.batch          batch             vm0          2  COMPLETED 
58.0           hostname         vm[0-1]          4  COMPLETED 
\end{lstlisting}

\textbf{Commentaires sur nodelist et alloccpus :}

\begin{table}[h]
\centering
\begin{tabular}{|l|l|c|p{7cm}|}
\hline
\textbf{Step} & \textbf{NodeList} & \textbf{AllocCPUS} & \textbf{Explication} \\
\hline
58 (job global) & vm[0-1] & 4 & Le job utilise 2 nœuds avec 4 CPUs au total \\
58.batch & vm0 & 2 & Le script batch s'exécute sur vm0 avec 2 CPUs \\
58.0 (hostname) & vm[0-1] & 4 & Le step utilise les 4 CPUs sur les 2 nœuds \\
\hline
\end{tabular}
\caption{Répartition des ressources par step}
\end{table}

\subsection{Q2 : Ajouter un second step avec sleep 120}

\textbf{Script \code{script-2.2.sh} :}
\begin{lstlisting}
#!/bin/bash
srun hostname
srun sleep 120
exit 0
\end{lstlisting}

\textbf{Suivi avec squeue :}
\begin{lstlisting}
[yahdhih.abdelwedoud@vm0 ~]$ squeue -s -j 61
         STEPID     NAME PARTITION     USER      TIME NODELIST
           61.1    sleep        vm yahdhih.      0:19 vm[0-1]
\end{lstlisting}

\textbf{Données d'accounting après terminaison :}
\begin{lstlisting}
[yahdhih.abdelwedoud@vm0 ~]$ sacct -j 61 -o jobid,jobname,nodelist,alloccpus,state,end,elapsed
       JobID    JobName        NodeList  AllocCPUS      State          End    Elapsed 
------------ ---------- --------------- ---------- ---------- ------------ ---------- 
61           script-2.+         vm[0-1]          4  COMPLETED 2026-02-02   00:02:01 
61.batch          batch             vm0          2  COMPLETED 2026-02-02   00:02:01 
61.0           hostname         vm[0-1]          4  COMPLETED 2026-02-02   00:00:01 
61.1              sleep         vm[0-1]          4  COMPLETED 2026-02-02   00:02:00 
\end{lstlisting}

\begin{table}[h]
\centering
\begin{tabular}{|l|c|p{8cm}|}
\hline
\textbf{Step} & \textbf{Elapsed} & \textbf{Explication} \\
\hline
61 (job global) & 00:02:01 & Durée totale du job \\
61.batch & 00:02:01 & Le batch attend la fin de tous les steps \\
61.0 (hostname) & 00:00:01 & Exécution instantanée \\
61.1 (sleep) & 00:02:00 & 120 secondes comme prévu \\
\hline
\end{tabular}
\caption{Durée d'exécution par step}
\end{table}

\subsection{Q3 : Boucle de 3 steps avec annulation (scancel)}

\textbf{Script \code{script-2.3.sh} :}
\begin{lstlisting}
#!/bin/bash
for i in 1 2 3; do
    srun sleep 120
done
exit 0
\end{lstlisting}

\textbf{Annulation du premier step :}
\begin{lstlisting}
[yahdhih.abdelwedoud@vm0 ~]$ scancel 62.0
\end{lstlisting}

\textbf{État final après annulation complète :}
\begin{lstlisting}
[yahdhih.abdelwedoud@vm0 ~]$ sacct -j 62 -o jobid,user,jobname%15,state%20,elapsed,exit
       JobID      User         JobName                State    Elapsed ExitCode 
------------ --------- --------------- -------------------- ---------- -------- 
62           yahdhih.+   script-2.3.sh    CANCELLED by 1003   00:05:02      0:0 
62.batch                         batch            CANCELLED   00:05:03     0:15 
62.0                             sleep    CANCELLED by 1003   00:01:40      0:9 
62.1                             sleep            COMPLETED   00:02:00      0:0 
62.2                             sleep            CANCELLED   00:01:22     0:15 
\end{lstlisting}

\textbf{Explication de la sortie de sacct :}

\begin{table}[h]
\centering
\begin{tabular}{|l|l|c|p{6cm}|}
\hline
\textbf{Step} & \textbf{State} & \textbf{ExitCode} & \textbf{Explication} \\
\hline
62 (job) & CANCELLED by 1003 & 0:0 & Job annulé par l'utilisateur (UID 1003) \\
62.batch & CANCELLED & 0:15 & Script batch interrompu (signal 15 = SIGTERM) \\
62.0 & CANCELLED by 1003 & 0:9 & Premier step annulé manuellement (signal 9 = SIGKILL) \\
62.1 & COMPLETED & 0:0 & Deuxième step terminé normalement \\
62.2 & CANCELLED & 0:15 & Troisième step interrompu \\
\hline
\end{tabular}
\caption{États des steps après annulation}
\end{table}

\textbf{Observations :}
\begin{itemize}
    \item \code{scancel 62.0} annule uniquement le step spécifié, le job continue avec les steps suivants
    \item \code{scancel 62} annule le job complet et tous les steps en cours
    \item L'ExitCode \code{0:15} indique une terminaison par signal SIGTERM (15)
    \item L'ExitCode \code{0:9} indique une terminaison par signal SIGKILL (9)
\end{itemize}

%%%%%%%%%%%%%%%%%%%%%%%%%%%%%%%%%%%%%%%%%%%%%%%%%%%%%%%%%%%%%%%%%%%%%%%%%%%%%%%
\section{Utilisation des QOS}
%%%%%%%%%%%%%%%%%%%%%%%%%%%%%%%%%%%%%%%%%%%%%%%%%%%%%%%%%%%%%%%%%%%%%%%%%%%%%%%

L'objectif est de comprendre le comportement des politiques d'ordonnancement en utilisant différentes QOS pour favoriser le démarrage de certains jobs.

\subsection{Q1 : Vérification du service slurmdbd et description des QOS}

\textbf{Vérification du service slurmdbd :}
\begin{lstlisting}
[yahdhih.abdelwedoud@vm0 ~]$ systemctl status slurmdbd
slurmdbd.service - Slurm DBD accounting daemon
   Loaded: loaded (/usr/lib/systemd/system/slurmdbd.service)
   Active: active (running) since lun. 2026-02-02 22:34:17 UTC
\end{lstlisting}

\textbf{Affichage des QOS disponibles :}
\begin{lstlisting}
[yahdhih.abdelwedoud@vm0 ~]$ sacctmgr show qos
      Name   Priority  GraceTime    Preempt PreemptMode  MaxWall MaxJobsPU
---------- ---------- ---------- ---------- ----------- -------- ---------
    normal          0   00:00:00            cluster                    
     debug         10   00:00:00            cluster     00:05:00        1
\end{lstlisting}

\textbf{Comparaison des deux QOS :}

\begin{table}[h]
\centering
\begin{tabular}{|l|c|c|}
\hline
\textbf{Propriété} & \textbf{normal} & \textbf{debug} \\
\hline
Priority & 0 & 10 \\
GraceTime & 00:00:00 & 00:00:00 \\
PreemptMode & cluster & cluster \\
UsageFactor & 1.0 & 1.0 \\
MaxWall & illimité & 00:05:00 (5 minutes) \\
MaxJobsPU & illimité & 1 \\
\hline
\end{tabular}
\caption{Comparaison des QOS normal et debug}
\end{table}

\textbf{Différences principales :}
\begin{enumerate}
    \item \textbf{Priority} : \code{debug} a une priorité de 10 contre 0 pour \code{normal}
    \item \textbf{MaxWall} : Les jobs \code{debug} sont limités à 5 minutes
    \item \textbf{MaxJobsPU} : Un seul job actif avec la QOS \code{debug} par utilisateur
\end{enumerate}

\subsection{Q2 : Lancement avec allocation exclusive et QOS debug}

\begin{lstlisting}
[yahdhih.abdelwedoud@vm0 ~]$ sbatch -n4 -N4 --exclusive script-2.2.sh
Submitted batch job 64
[yahdhih.abdelwedoud@vm0 ~]$ sbatch -n4 -N4 --exclusive script-2.2.sh
Submitted batch job 65
[yahdhih.abdelwedoud@vm0 ~]$ sbatch -n4 -N4 --exclusive --qos debug script-2.2.sh
Submitted batch job 66
\end{lstlisting}

La QOS \code{debug} ayant une priorité plus élevée (10 vs 0), les jobs debug seront favorisés dès que les ressources se libèrent.

\subsection{Q3 : Lancement sans allocation exclusive}

\begin{lstlisting}
[yahdhih.abdelwedoud@vm0 ~]$ squeue -O jobid,state,qos,timeused,reason
JOBID     STATE     QOS      TIME     REASON              
64        PENDING   normal   0:00     Resources           
65        PENDING   normal   0:00     Resources           
68        PENDING   normal   0:00     Priority            
69        PENDING   normal   0:00     Priority            
63        RUNNING   normal   4:53     None                
\end{lstlisting}

\begin{table}[h]
\centering
\begin{tabular}{|l|p{9cm}|}
\hline
\textbf{REASON} & \textbf{Signification} \\
\hline
Resources & Le job attend que des ressources se libèrent \\
Priority & Le job attend car d'autres jobs ont une priorité plus élevée \\
None & Le job est en cours d'exécution \\
\hline
\end{tabular}
\caption{Signification des raisons d'attente}
\end{table}

\subsection{Q4 : Lancement avec QOS debug}

\begin{lstlisting}
[yahdhih.abdelwedoud@vm0 ~]$ sbatch -n4 -N4 --qos debug script-2.2.sh
Submitted batch job 71
[yahdhih.abdelwedoud@vm0 ~]$ sbatch -n4 -N4 --qos debug script-2.2.sh
Submitted batch job 72
\end{lstlisting}

\textbf{Constatations :}
\begin{enumerate}
    \item Les jobs debug apparaissent en tête de la file d'attente grâce à leur priorité plus élevée
    \item Le job 72 attend avec REASON = Priority à cause de MaxJobsPU = 1
    \item Le job 71 attend les Resources et sera le prochain à démarrer
\end{enumerate}

\textbf{Conclusion :} La QOS \code{debug} est conçue pour des tests rapides :
\begin{itemize}
    \item \textbf{Avantage :} Priorité élevée pour un démarrage rapide
    \item \textbf{Contraintes :} Durée limitée (5 min) et un seul job à la fois par utilisateur
\end{itemize}

%%%%%%%%%%%%%%%%%%%%%%%%%%%%%%%%%%%%%%%%%%%%%%%%%%%%%%%%%%%%%%%%%%%%%%%%%%%%%%%
\section{Comportements aux limites}
%%%%%%%%%%%%%%%%%%%%%%%%%%%%%%%%%%%%%%%%%%%%%%%%%%%%%%%%%%%%%%%%%%%%%%%%%%%%%%%

L'idée est de constater le comportement de Slurm lorsqu'un job batch ne se comporte pas de la façon attendue ou que le système lui-même introduit des problèmes.

\subsection{Q1 : Limite de temps d'exécution (time limit)}

\textbf{Test avec un script dépassant la limite de temps :}
\begin{lstlisting}
[yahdhih.abdelwedoud@vm0 ~]$ sbatch -t 1 script-2.3.sh
Submitted batch job 74
[yahdhih.abdelwedoud@vm0 ~]$ sacct -j 74 -o jobid,user,jobname%15,state,elapsed
       JobID      User         JobName      State    Elapsed 
------------ --------- --------------- ---------- ---------- 
74           yahdhih.+   script-2.3.sh    TIMEOUT   00:01:04 
74.batch                         batch  CANCELLED   00:01:05 
74.0                             sleep  CANCELLED   00:01:04 
\end{lstlisting}

\begin{table}[h]
\centering
\begin{tabular}{|l|l|c|p{6cm}|}
\hline
\textbf{Step} & \textbf{State} & \textbf{Elapsed} & \textbf{Explication} \\
\hline
74 (job) & TIMEOUT & 00:01:04 & Le job a dépassé la limite de temps \\
74.batch & CANCELLED & 00:01:05 & Le batch a été annulé suite au timeout \\
74.0 (sleep) & CANCELLED & 00:01:04 & Le step en cours a été annulé \\
\hline
\end{tabular}
\caption{États après dépassement du temps limite}
\end{table}

\textbf{Réponses :}
\begin{itemize}
    \item \textbf{State associé au job :} \code{TIMEOUT}
    \item \textbf{Nombre de steps exécutés :} 1 seul step a été lancé avant le timeout
\end{itemize}

\subsection{Q2 : Limite de mémoire par CPU}

\textbf{Script \code{gen100Mo.sh} :}
\begin{lstlisting}
#!/bin/bash
srun dd if=/dev/zero of=/dev/shm/128Mo bs=1M count=128
sleep 30
rm /dev/shm/128Mo
exit 0
\end{lstlisting}

\textbf{Lancement avec limite de mémoire :}
\begin{lstlisting}
[yahdhih.abdelwedoud@vm0 ~]$ sbatch --mem-per-cpu=100 -n1 gen100Mo.sh
Submitted batch job 76
[yahdhih.abdelwedoud@vm0 ~]$ sacct -j 76 -o jobid,user,jobname%15,state,elapsed,exit
       JobID      User         JobName      State    Elapsed ExitCode 
------------ --------- --------------- ---------- ---------- -------- 
76           yahdhih.+     gen100Mo.sh  COMPLETED   00:00:30      0:0 
76.batch                         batch  COMPLETED   00:00:30      0:0 
76.0                                dd  CANCELLED   00:00:00      0:9 
\end{lstlisting}

\textbf{Explication :}
\begin{enumerate}
    \item Le step \code{dd} a été annulé (CANCELLED) avec ExitCode \code{0:9} car il a tenté d'allouer 128Mo alors que la limite était de 100Mo
    \item Le signal 9 (SIGKILL) indique que Slurm a forcé la terminaison du processus
\end{enumerate}

\subsection{Q3 : Défaillance d'un nœud (NODE\_FAIL)}

\textbf{Reboot du nœud vm1 pendant l'exécution :}
\begin{lstlisting}
[yahdhih.abdelwedoud@vm0 ~]$ sudo ssh vm1 reboot
Connection to vm1 closed by remote host.
\end{lstlisting}

\textbf{État final du job :}
\begin{lstlisting}
[yahdhih.abdelwedoud@vm0 ~]$ sacct -j 77 -o jobid,user,jobname%15,state,elapsed,exit
       JobID      User         JobName      State    Elapsed ExitCode 
------------ --------- --------------- ---------- ---------- -------- 
77           yahdhih.+   script-2.3.sh  NODE_FAIL   00:03:48      1:0 
77.batch                         batch  CANCELLED   00:03:49     0:15 
77.0                             sleep  COMPLETED   00:01:30      0:0 
77.1                             sleep  COMPLETED   00:01:30      0:0 
77.2                             sleep  NODE_FAIL   00:00:47      0:0 
\end{lstlisting}

\textbf{Conclusion :} Slurm gère les défaillances de nœuds en marquant les jobs affectés avec l'état \code{NODE\_FAIL} et en mettant les nœuds défaillants en état \code{down}.

%%%%%%%%%%%%%%%%%%%%%%%%%%%%%%%%%%%%%%%%%%%%%%%%%%%%%%%%%%%%%%%%%%%%%%%%%%%%%%%
\section{Gestion des nœuds}
%%%%%%%%%%%%%%%%%%%%%%%%%%%%%%%%%%%%%%%%%%%%%%%%%%%%%%%%%%%%%%%%%%%%%%%%%%%%%%%

L'état actuel du cluster n'est pas complètement opérationnel. Le reboot de vm1 a changé son état et ne le rend plus éligible à l'allocation.

\subsection{Q1 : Lancement d'un script nécessitant 4 nœuds exclusivement}

\begin{lstlisting}
[yahdhih.abdelwedoud@vm0 ~]$ sbatch -n4 -N4 --exclusive script-2.3.sh
Submitted batch job 58
[yahdhih.abdelwedoud@vm0 ~]$ squeue
    JOBID PARTITION   NAME     USER ST   TIME  NODES NODELIST(REASON)
       58        vm script-2 yahdhih. PD   0:00      4 (Resources)
\end{lstlisting}

\textbf{Constatation :} Le job est en attente (\textbf{PENDING}) avec la raison \textbf{Resources} car vm1 est down.

\subsection{Q2 : Remise en production du nœud vm1}

\begin{lstlisting}
[yahdhih.abdelwedoud@vm0 ~]$ sudo scontrol update nodename=vm1 state=idle
[yahdhih.abdelwedoud@vm0 ~]$ sinfo
PARTITION AVAIL  TIMELIMIT  NODES  STATE NODELIST
vm*          up   infinite      4  alloc vm[0-3]
\end{lstlisting}

Le job 58 passe automatiquement en exécution (\textbf{RUNNING}) sur les 4 nœuds.

\subsection{Q3 : Utilisation de l'option --requeue}

\begin{table}[h]
\centering
\begin{tabular}{|l|p{5cm}|p{5cm}|}
\hline
\textbf{Aspect} & \textbf{Sans --requeue} & \textbf{Avec --requeue} \\
\hline
État après NODE\_FAIL & Job terminé avec NODE\_FAIL & Job remis en file d'attente (PENDING) \\
Après remise en prod & Job reste NODE\_FAIL & Job reprend l'exécution \\
Comportement & Le job est définitivement échoué & Le job est automatiquement relancé \\
\hline
\end{tabular}
\caption{Comparaison avec et sans --requeue}
\end{table}

\subsection{Q4 : Exclusivité et nœud drainé}

Un nœud drainé finit les tâches qui lui sont actuellement attribuées, puis se met hors service.

\begin{lstlisting}
[yahdhih.abdelwedoud@vm0 ~]$ sudo scontrol update node=vm1 state=drain reason="drainage"
[yahdhih.abdelwedoud@vm0 ~]$ sinfo
PARTITION AVAIL  TIMELIMIT  NODES  STATE NODELIST
vm*          up   infinite      1   drng vm1
vm*          up   infinite      3  alloc vm[0,2-3]
\end{lstlisting}

Le nœud vm1 passe de "drng" (draining) à "drain" (drainé) une fois le job terminé.

\subsection{Q5 : Passage d'un nœud en état "down"}

\begin{lstlisting}
[yahdhih.abdelwedoud@vm0 ~]$ sudo scontrol update node=vm1 state=down reason="down"
[yahdhih.abdelwedoud@vm0 ~]$ sacct -j 62
       JobID    JobName  Partition    Account  AllocCPUS      State ExitCode 
------------ ---------- ---------- ---------- ---------- ---------- -------- 
62           script-2.+         vm     guests          4  NODE_FAIL      1:0 
\end{lstlisting}

Le job 62 est immédiatement terminé avec l'état \textbf{NODE\_FAIL}.

\subsection{Q6 : Différence entre "drain" et "down"}

\begin{table}[h]
\centering
\begin{tabular}{|l|p{5cm}|p{5cm}|}
\hline
\textbf{Aspect} & \textbf{drain} & \textbf{down} \\
\hline
Jobs en cours & Continuent jusqu'à leur fin & Interrompus immédiatement \\
Nouveaux jobs & Non acceptés & Non acceptés \\
État du job affecté & COMPLETED (termine normalement) & NODE\_FAIL (échec) \\
Transition d'état & idle → drng → drain & idle → down \\
Usage typique & Maintenance planifiée & Panne ou urgence \\
\hline
\end{tabular}
\caption{Comparaison drain vs down}
\end{table}

\textbf{Pour remettre un nœud en production :}
\begin{lstlisting}
sudo scontrol update node=<nom_noeud> state=resume
# ou
sudo scontrol update nodename=<nom_noeud> state=idle
\end{lstlisting}

%%%%%%%%%%%%%%%%%%%%%%%%%%%%%%%%%%%%%%%%%%%%%%%%%%%%%%%%%%%%%%%%%%%%%%%%%%%%%%%
\section{Test du Backfilling}
%%%%%%%%%%%%%%%%%%%%%%%%%%%%%%%%%%%%%%%%%%%%%%%%%%%%%%%%%%%%%%%%%%%%%%%%%%%%%%%

Le cluster est pleinement opérationnel. Cette partie illustre le fonctionnement du « backfilling » - une technique d'ordonnancement où des jobs de moindre priorité peuvent passer avant des jobs prioritaires s'ils peuvent se terminer avant que les ressources ne soient disponibles pour le job prioritaire.

\textbf{Scripts utilisés :}
\begin{itemize}
    \item \textbf{bkf-1.sh} : Utilise une partie des ressources (2 nœuds) pour une période longue
    \item \textbf{bkf-2.sh} : Nécessite l'intégralité des ressources (4 nœuds exclusifs) avec grande priorité
    \item \textbf{bkf-3.sh} : Moins prioritaire, durée longue (non backfillable)
    \item \textbf{bkf-4.sh} : Moins prioritaire, durée courte (backfillable)
\end{itemize}

\subsection{Q1 : Mise en place du scénario de backfilling}

\textbf{Lancement des 4 jobs :}
\begin{lstlisting}
[yahdhih.abdelwedoud@vm0 ~]$ sbatch -t 2 -n2 -N2 bkf-1.sh ; \
  sbatch -t 3 -n4 -N4 --exclusive --nice=50 bkf-2.sh ; \
  sbatch -t 4 -n2 -N2 --nice=100 bkf-3.sh ; \
  sbatch -t 1 -n1 -N1 --nice=100 bkf-4.sh 
Submitted batch job 69
Submitted batch job 70
Submitted batch job 71
Submitted batch job 72
\end{lstlisting}

\textbf{Configuration des jobs :}

\begin{table}[h]
\centering
\begin{tabular}{|c|l|c|c|c|l|}
\hline
\textbf{Job ID} & \textbf{Script} & \textbf{Nœuds} & \textbf{Temps max} & \textbf{Nice} & \textbf{Priorité} \\
\hline
69 & bkf-1.sh & 2 (-N2) & 2 min & 0 (défaut) & Haute \\
70 & bkf-2.sh & 4 (--exclusive) & 3 min & 50 & Moyenne \\
71 & bkf-3.sh & 2 (-N2) & 4 min & 100 & Basse \\
72 & bkf-4.sh & 1 (-N1) & 1 min & 100 & Basse \\
\hline
\end{tabular}
\caption{Configuration des jobs pour le backfilling}
\end{table}

\textbf{Note :} Plus la valeur de \code{--nice} est élevée, plus la priorité est basse.

\textbf{Observation du backfilling :}
\begin{lstlisting}
[yahdhih.abdelwedoud@vm0 ~]$ squeue -O jobid,qos,timeused,state,reason,nodelist,priority
JOBID   QOS      TIME   STATE     REASON     NODELIST   PRIORITY            
70      normal   0:00   PENDING   Resources             0.00001782574691    
71      normal   0:08   RUNNING   None       vm[2-3]    0.00001781410538    
\end{lstlisting}

\textbf{Le job 71 (priorité basse) s'exécute AVANT le job 70 (priorité plus haute) !}

C'est le backfilling en action :
\begin{enumerate}
    \item Le job 70 nécessite 4 nœuds exclusivement → il doit attendre que TOUS les nœuds soient libres
    \item Le job 71 ne nécessite que 2 nœuds → il peut utiliser les nœuds libérés par d'autres jobs
    \item Slurm calcule que le job 71 peut se terminer AVANT que les 4 nœuds ne soient disponibles pour le job 70
    \item Donc le job 71 est « backfillé » → il passe avant le job 70 sans retarder ce dernier
\end{enumerate}

\subsection{Q2 : Annulation du premier job - Limitation du backfilling}

\begin{lstlisting}
[yahdhih.abdelwedoud@vm0 ~]$ scancel 73
[yahdhih.abdelwedoud@vm0 ~]$ squeue -O jobid,qos,timeused,state,reason,nodelist,priority
JOBID   QOS      TIME   STATE     REASON     NODELIST   PRIORITY            
74      normal   0:00   PENDING   Resources             0.00001782481559    
75      normal   0:35   RUNNING   None       vm[2-3]    0.00001781317406    
76      normal   0:35   RUNNING   None       vm0        0.00001781317406    
\end{lstlisting}

\textbf{Observation critique :} Le job 74 (priorité plus haute) est \textbf{toujours en attente} alors que les jobs 75 et 76 (priorité plus basse) continuent de s'exécuter !

\textbf{Conclusion sur les limitations du backfilling :}
\begin{enumerate}
    \item \textbf{Pas de préemption :} Une fois qu'un job est backfillé et en cours d'exécution, il ne peut pas être interrompu
    \item \textbf{Décision irrévocable :} Les décisions de backfilling sont prises au moment de l'ordonnancement
    \item \textbf{Inversion de priorité :} En cas de changement de contexte, des jobs de basse priorité peuvent bloquer des jobs de haute priorité
\end{enumerate}

\subsection{Q3 : Sans temps d'allocation - Besoins du backfilling}

\textbf{Lancement sans limites de temps :}
\begin{lstlisting}
[yahdhih.abdelwedoud@vm0 ~]$ sbatch -n2 -N2 bkf-1.sh ; \
  sbatch -n4 -N4 --exclusive bkf-2.sh ; \
  sbatch -n2 -N2 bkf-3.sh ; \
  sbatch -n1 -N1 bkf-4.sh 
Submitted batch job 77
Submitted batch job 78
Submitted batch job 79
Submitted batch job 80
\end{lstlisting}

\textbf{Note :} Pas d'option \code{-t} (temps), pas de \code{--nice} (même priorité pour tous).

\begin{lstlisting}
[yahdhih.abdelwedoud@vm0 ~]$ squeue -O jobid,qos,timeused,state,reason,nodelist,priority
JOBID   QOS      TIME   STATE     REASON     NODELIST   PRIORITY            
78      normal   0:00   PENDING   Resources             0.00001783645712    
77      normal   1:47   RUNNING   None       vm[2-3]    0.00001783645712    
79      normal   1:47   RUNNING   None       vm[2-3]    0.00001783645712    
\end{lstlisting}

\textbf{Analyse : Pourquoi pas de backfilling ?}

\begin{table}[h]
\centering
\begin{tabular}{|p{6cm}|p{6cm}|}
\hline
\textbf{Avec -t (temps spécifié)} & \textbf{Sans -t} \\
\hline
Slurm connaît la durée maximale des jobs & Durée = UNLIMITED \\
Peut calculer quand les ressources seront libres & Impossible de prédire la fin des jobs \\
Peut décider si un job peut être backfillé & Aucun backfilling possible \\
\hline
\end{tabular}
\caption{Impact du temps sur le backfilling}
\end{table}

\textbf{Besoins du backfilling :}
\begin{enumerate}
    \item \textbf{TimeLimit défini :} L'option \code{-t} ou \code{--time} doit être spécifiée
    \item \textbf{Différences de priorité :} Les jobs doivent avoir des priorités différentes
    \item \textbf{Ressources partielles :} Le job à backfiller doit nécessiter moins de ressources
\end{enumerate}

\textbf{Conclusion :} Sans spécification de temps (\code{-t}), Slurm ne peut pas prédire quand les jobs en cours se termineront, donc le backfilling est \textbf{désactivé} par sécurité.

\textbf{Recommandation :} Toujours spécifier un temps maximum avec \code{-t} ou \code{--time} pour permettre un ordonnancement optimal avec backfilling.

\end{document}
